\begin{table*}[hbtp]
\centering
\footnotesize
\setlength{\tabcolsep}{3.0pt}
\renewcommand{\arraystretch}{1.09}
\caption{CIFAR-100-C under continual adaptation over the 15-corruption stream, with all methods using the same ResNeXt-29 backbone. Results are reported as mean $\pm$ std over three random orders. Best in \textbf{bold}; second-best \underline{underlined}.}
\label{tab:l1a}
\resizebox{\textwidth}{!}{%
\begin{tabular}{lccccccccccccccc!{\color{black!35}\vrule width 0.45pt}c}
\toprule
Method & Bright. & Contr. & Defoc. & Elastic & Fog & Frost & Gauss. & Glass & Impulse & JPEG & Motion & Pixel. & Shot & Snow & Zoom & \avg \\
\midrule
Source & \res{70.5}{0.0} & \res{44.9}{0.0} & \res{70.7}{0.0} & \res{62.8}{0.0} & \res{49.7}{0.0} & \res{54.2}{0.0} & \res{27.0}{0.0} & \res{45.9}{0.0} & \res{60.6}{0.0} & \res{58.8}{0.0} & \res{69.2}{0.0} & \res{25.3}{0.0} & \res{32.0}{0.0} & \res{60.5}{0.0} & \res{71.2}{0.0} & \res{53.6}{0.0} \\
Tent (ICLR'21) & \res{74.5}{0.3} & \res{67.4}{3.7} & \res{71.3}{2.6} & \res{66.0}{1.2} & \res{62.8}{1.0} & \res{64.9}{1.5} & \res{62.2}{1.3} & \res{62.4}{1.2} & \res{58.9}{0.6} & \res{60.2}{2.4} & \res{68.6}{2.7} & \res{68.2}{2.8} & \res{63.7}{2.8} & \res{65.2}{2.7} & \res{68.4}{0.8} & \res{65.6}{0.4} \\
EATA (ICML'22) & \res{74.0}{0.3} & \res{70.2}{0.5} & \res{72.5}{0.6} & \res{66.7}{0.3} & \res{63.1}{1.0} & \res{67.7}{0.7} & \res{62.3}{0.2} & \res{62.4}{0.1} & \res{61.9}{0.2} & \res{60.7}{0.4} & \res{70.1}{0.6} & \res{69.8}{0.4} & \res{64.7}{0.3} & \res{67.1}{0.7} & \res{72.5}{0.5} & \res{67.1}{0.0} \\
SAR (ICLR'23) & \res{73.3}{0.1} & \res{69.7}{0.0} & \res{72.4}{0.0} & \res{66.1}{1.2} & \res{61.7}{0.7} & \res{66.5}{0.2} & \res{61.2}{0.9} & \res{61.1}{1.4} & \best{63.3}{1.0} & \res{61.0}{1.1} & \res{70.6}{0.1} & \res{68.1}{0.0} & \res{62.9}{1.3} & \res{68.0}{0.1} & \res{71.9}{0.1} & \res{66.5}{0.1} \\
CoTTA (CVPR'22) & \res{72.0}{0.0} & \res{65.9}{0.5} & \res{69.6}{0.3} & \res{65.7}{0.3} & \res{58.7}{0.2} & \res{66.8}{0.1} & \res{62.4}{0.3} & \ulres{62.9}{0.2} & \res{61.8}{0.0} & \best{63.6}{0.9} & \res{69.0}{0.4} & \res{69.8}{0.1} & \res{63.6}{0.6} & \res{65.4}{0.2} & \res{69.7}{0.1} & \res{65.8}{0.1} \\
LCoTTA (NeurIPS'25) & \res{73.9}{0.3} & \res{70.4}{0.8} & \res{72.6}{0.7} & \res{66.7}{0.3} & \res{62.6}{1.1} & \res{67.9}{0.6} & \res{62.4}{0.2} & \res{62.6}{0.6} & \res{61.3}{1.0} & \res{60.7}{0.3} & \res{70.3}{0.4} & \res{69.9}{0.8} & \res{64.6}{0.7} & \res{67.1}{0.6} & \res{72.8}{0.2} & \res{67.0}{0.0} \\
AdaDEM (NeurIPS'25) & \res{73.9}{0.4} & \res{71.3}{0.7} & \res{74.0}{0.6} & \res{66.7}{0.8} & \res{61.4}{0.3} & \res{68.2}{0.6} & \res{61.3}{1.0} & \res{61.6}{0.3} & \res{62.0}{0.1} & \res{59.9}{1.1} & \ulres{72.0}{0.5} & \res{69.8}{1.1} & \res{63.1}{1.6} & \res{67.5}{0.3} & \best{74.6}{0.0} & \res{67.1}{0.0} \\
SURGEON (CVPR'25) & \ulres{74.7}{1.4} & \res{65.8}{6.3} & \res{68.3}{4.5} & \res{65.9}{1.9} & \res{62.6}{2.7} & \res{64.7}{3.5} & \ulres{63.8}{1.6} & \res{62.9}{1.0} & \res{57.4}{1.1} & \ulres{63.5}{1.4} & \res{66.2}{3.9} & \res{68.8}{1.9} & \ulres{65.6}{1.0} & \res{66.0}{1.2} & \res{65.5}{2.7} & \res{65.4}{1.3} \\
DeYO (ICLR'24) & \best{74.8}{0.1} & \ulres{72.1}{0.4} & \ulres{74.2}{0.4} & \ulres{67.7}{0.2} & \ulres{63.6}{0.6} & \ulres{68.7}{0.1} & \res{62.8}{0.2} & \res{62.4}{0.2} & \res{62.4}{0.2} & \res{62.3}{0.6} & \res{71.8}{0.3} & \ulres{70.7}{0.2} & \res{64.7}{0.3} & \ulres{68.1}{0.2} & \res{74.0}{0.2} & \ulres{68.0}{0.1} \\
RoTTA (CVPR'23) & \res{71.6}{0.9} & \res{60.6}{12.8} & \res{72.5}{2.1} & \res{66.2}{1.1} & \res{61.7}{0.3} & \res{68.5}{0.8} & \res{58.6}{3.0} & \res{59.4}{2.0} & \res{59.2}{0.3} & \res{57.0}{5.0} & \res{71.1}{1.4} & \res{67.1}{4.0} & \res{59.4}{4.0} & \res{67.4}{1.3} & \ulres{74.2}{0.5} & \res{65.0}{0.1} \\
\midrule
\rowcolor{gray!10}
\textbf{\method{}} & \res{74.6}{0.6} & \best{72.8}{0.3} & \best{74.4}{0.6} & \best{68.5}{0.1} & \best{64.7}{0.6} & \best{70.2}{0.5} & \best{64.3}{0.2} & \best{64.9}{0.1} & \ulres{63.2}{0.3} & \res{61.8}{1.1} & \best{72.2}{0.4} & \best{71.4}{0.5} & \best{66.2}{0.5} & \best{69.1}{0.3} & \res{74.0}{0.0} & \best{68.8}{0.2} \\
\bottomrule
\end{tabular}}
\end{table*}

\section{Experiments and Results}

\subsection{Experimental setup and five-level protocol}

The evaluation follows a five-level protocol. Level 1 studies CIFAR-100-C and ImageNet-C under standard TTA, where the model resets between corruptions, and continual TTA, where one model processes the full stream without gradient reset~\citep{wang2021tent,wang2022cotta}. Both use the $15$ ImageNet-C corruption types at severity $5$~\citep{hendrycks2019imagenetc}. Level 2 tests natural shifts on ImageNet-A, ImageNet-R, ImageNet-V2, and ImageNet-Sketch with ViT-B/16~\citep{dosovitskiy2021vit,hendrycks2021imageneta,hendrycks2021manyfaces,recht2021imagenetv2}. Level 3 evaluates wild non-i.i.d. streams under label shift, mixed shifts, and batch-size-one inference, following WILDS~\citep{koh2021wilds}. Level 4 studies Cityscapes$\rightarrow$ACDC segmentation with SegFormer-B5~\citep{cordts2016cityscapes,sakaridis2021acdc,xie2021segformer}. Level 5 studies prompt-only adaptation for CLIP ViT-B/16~\citep{radford2021clip}. Classification and VLM experiments report top-1 accuracy; segmentation reports online error and mIoU. Means $\pm$ standard deviations are reported when available. Standard corruption runs use three random seeds; continual runs use three random corruption orders. The protocol isolates unreliable evidence, online drift, and entity-scale mismatch.

\subsection{Level 1: Synthetic corruptions}


\begin{wraptable}{r}{0.60\columnwidth}
\vspace{-1.2\baselineskip}
\centering
\normalsize
\setlength{\tabcolsep}{3pt}
\renewcommand{\arraystretch}{1.05}
\caption{\textbf{ImageNet-C under standard and continual adaptation.}
Standard columns report mean $\pm$ std across three random seeds; continual columns report mean $\pm$ std across three random corruption orders. Unsupported entries are excluded from best-baseline margins; collapse entries are raw official-scope reproductions and are not used as the primary comparison. Best in \textbf{bold}; second-best \underline{underlined}.}
\label{tab:l2}
\resizebox{\linewidth}{!}{%
\begin{tabular}{@{}lcccc@{}}
\toprule
& \multicolumn{2}{c}{ResNet-50-GN} & \multicolumn{2}{c}{ViT-B/16} \\
\cmidrule(lr){2-3}\cmidrule(lr){4-5}
Method & Standard & Continual & Standard & Continual \\
\midrule
Source & \res{31.44}{0.00} & \res{31.44}{0.00} & \res{39.83}{0.00} & \res{39.83}{0.00} \\
Tent (ICLR'21) & \res{32.02}{0.00} & \res{21.84}{1.38} & \res{54.62}{0.03} & \res{49.02}{0.14} \\
EATA (ICML'22) & \res{41.89}{0.02} & \ulres{34.82}{0.57} & \res{58.25}{1.09} & \res{18.04}{8.68} \\
SAR (ICLR'23) & \res{35.84}{0.00} & \res{32.30}{8.20} & \res{54.57}{0.03} & \res{55.88}{0.58} \\
CoTTA (CVPR'22) & \res{31.49}{0.00} & \res{31.44}{0.03} & \res{40.25}{0.00} & \res{42.10}{0.09} \\
Lifelong TTA (NeurIPS'25) & \res{40.36}{0.00} & \res{34.02}{0.79} & \ulres{60.18}{0.10} & \res{36.73}{0.20} \\
AdaDEM (NeurIPS'25) & \res{32.21}{0.00} & \res{30.72}{0.29} & \res{58.94}{1.83} & \ulres{64.11}{0.43} \\
FOA (ICML'24) & \multicolumn{1}{c}{\pending} & \multicolumn{1}{c}{\pending} & \res{48.98}{0.18} & \res{48.56}{0.34} \\
DeYO (ICLR'24) & \ulres{43.79}{0.01} & \res{23.98}{13.91} & \res{60.00}{0.59} & \res{56.52}{0.40} \\
\midrule
\rowcolor{gray!8}
\textbf{\method{}} & \best{48.40}{0.07} & \best{45.84}{0.08} & \best{60.21}{0.03} & \best{65.14}{0.05} \\
\bottomrule
\end{tabular}}
\vspace{-0.8\baselineskip}
\end{wraptable}

% \begin{table*}[t]
% \centering
% \small
% \setlength{\tabcolsep}{6pt}
% \renewcommand{\arraystretch}{1.08}
% \caption{\textbf{ImageNet-C under standard and continual adaptation.}
% Results are reported as mean \(\pm\) std across three random orders. The best result in each column is highlighted in \textbf{bold}, and the second-best is \underline{underlined}.}
% \label{tab:l2}
% \begin{tabular}{@{}lcccc@{}}
% \toprule
% & \multicolumn{2}{c}{ResNet-50-GN} & \multicolumn{2}{c}{ViT-B/16} \\
% \cmidrule(lr){2-3}\cmidrule(lr){4-5}
% Method & Standard & Continual & Standard & Continual \\
% \midrule
% Source & \res{31.44}{0.00} & \res{31.44}{0.00} & \res{39.83}{0.00} & \res{39.83}{0.00} \\
% Tent (ICLR'21) & \res{32.02}{0.00} & \res{21.84}{1.38} & \res{54.62}{0.03} & \res{49.02}{13.95} \\
% EATA (ICML'22) & \res{41.89}{0.02} & \ulres{34.82}{0.57} & \res{58.25}{1.09} & \res{18.04}{8.68} \\
% SAR (ICLR'23) & \res{35.84}{0.00} & \res{32.30}{8.20} & \res{54.57}{0.03} & \res{55.88}{0.58} \\
% CoTTA (CVPR'22) & \res{31.49}{0.00} & \res{31.44}{0.03} & \res{40.25}{0.00} & \res{42.10}{0.09} \\
% Lifelong TTA (NeurIPS'25) & \res{40.36}{0.00} & \res{34.02}{0.79} & \ulres{60.18}{0.10} & \res{36.73}{0.20} \\
% AdaDEM (NeurIPS'25) & \res{32.21}{0.00} & \res{30.72}{0.29} & \res{58.94}{1.83} & \ulres{64.11}{0.43} \\
% SURGEON (CVPR'25) & \res{0.18}{0.01} & \res{0.13}{0.01} & \res{5.03}{0.57} & \res{3.17}{2.05} \\
% FOA (ICML'24) & \multicolumn{1}{c}{\pending} & \multicolumn{1}{c}{\pending} & \res{48.98}{0.18} & \res{48.56}{0.34} \\
% DeYO (ICLR'24) & \ulres{43.79}{0.01} & \res{23.98}{13.91} & \res{60.00}{0.59} & \res{56.52}{0.40} \\
% \midrule
% \rowcolor{gray!8}
% \textbf{\method{}} & \best{48.40}{0.07} & \best{45.84}{0.08} & \best{60.21}{0.03} & \best{65.14}{0.05} \\
% \bottomrule
% \end{tabular}
% \end{table*}
% % \begin{table}[htbp]
% % \centering
% % \small
% % \caption{ImageNet-C under standard and continual adaptation. Entries marked with - denote absent or invalid reproductions in the current artifact set.}
% % \label{tab:l2}
% % \resizebox{\textwidth}{!}{%
% % \begin{tabular}{llcc|llcc}
% % \toprule
% % Backbone & Method & Standard & Continual & Backbone & Method & Standard & Continual \\
% % \midrule
% % \multirow{12}{*}{ResNet-50-GN}
% % & Source & 31.44 & 31.44
% % & \multirow{12}{*}{ViT-B/16}
% % & Source & 39.83 & 39.83 \\
% % & Tent (ICLR'21) & 32.02 & 21.84
% % && Tent (ICLR'21) & 54.62 & 49.02 \\
% % & EATA (ICML'22) & 41.89 & 34.82
% % && EATA (ICML'22) & 58.25 & 18.04 \\
% % & SAR (ICLR'23) & 35.84 & 32.30
% % && SAR (ICLR'23) & 54.57 & 55.88 \\
% % & CoTTA (CVPR'22) & 31.49 & 31.44
% % && CoTTA (CVPR'22) & 40.25 & 42.10 \\
% % & Lifelong TTA (NeurIPS'25) & 40.36 & 34.02
% % && Lifelong TTA (NeurIPS'25) & 60.18 & 36.73 \\
% % & AdaDEM (NeurIPS'25) & 32.21 & 30.72
% % && AdaDEM (NeurIPS'25) & 58.94 & 64.11 \\
% % & SURGEON (CVPR'25) & 0.18 & 0.13
% % && SURGEON (CVPR'25) & 5.03 & 3.17 \\
% % & FOA (ICML'24) & - & -
% % && FOA (ICML'24) & 48.98 & 48.56 \\
% % & DeYO (ICLR'24) & 43.79 & 23.98
% % && DeYO (ICLR'24) & 60.00 & 56.52 \\
% % % & RoTTA (CVPR'23) & 31.44 & 31.44
% % % && RoTTA (CVPR'23) & 39.83 & 39.83 \\
% % & \method{} & 31.68 & 31.95
% % && \method{} & 41.41 & 42.43 \\
% % \bottomrule
% % \end{tabular}}
% % \end{table}


Synthetic corruptions test whether a TTA method can reject unsafe confidence while still adapting to degraded inputs. On CIFAR-100-C continual adaptation, Table~\ref{tab:l1a} shows \method{} at $68.8\pm0.2$, beating DeYO by $0.8$ points and EATA/AdaDEM by $1.7$. It is best on $11$ of $15$ corruptions, covering both photometric and structure-changing cases. Appendix~\ref{app:additional_cifar} reports the standard CIFAR-100-C protocol, where \method{} reaches $69.1\pm0.0$, $0.4$ points above SURGEON become new best method.

ImageNet-C exposes the same reliability issue across backbones. In Table~\ref{tab:l2}, \method{} reaches $48.40\pm0.07$ on ResNet-50-GN standard adaptation and $45.84\pm0.08$ on ResNet-50-GN continual adaptation, improving over the strongest baselines by $4.61$ and $11.02$ points. On ViT-B/16, it reaches $60.21\pm0.03$ in the standard setting and $65.14\pm0.05$ in the continual setting, improving over the best baselines by $0.03$ and $1.03$ points. The small ViT standard margin and larger continual margins suggest that \method{} helps most when stream order and drift matter.

\subsection{Level 2: Natural distribution shifts}
\begin{table*}[b]
\centering
\begin{minipage}[t]{0.59\textwidth}
\centering
\footnotesize
\setlength{\tabcolsep}{3.5pt}
\captionof{table}{Natural-shift OOD performance on ViT-B/16. Collapse entries are raw official-scope reproductions and are not used as the primary comparison. Best in bold; second-best underlined.}
\label{tab:l3}

\adjustbox{max width=\linewidth}{%
\begin{tabular}{lccccc}
\toprule
Method & IN-A & IN-R & IN-V2 & IN-Sketch & Avg. \\
\midrule
Source   & \res{15.61}{0.00} & \res{35.14}{0.00} & \res{69.42}{0.00} & \res{27.85}{0.00} & \res{37.01}{0.00} \\
Tent     & \res{17.95}{0.02} & \res{38.23}{0.05} & \ulres{69.77}{0.17} & \res{33.72}{2.09} & \res{39.92}{0.58} \\
EATA     & \res{17.04}{0.02} & \res{35.75}{0.80} & \res{69.62}{0.07} & \res{38.99}{0.28} & \res{40.35}{0.21} \\
SAR      & \res{17.78}{0.06} & \res{36.61}{0.05} & \res{69.56}{0.07} & \res{33.91}{0.06} & \res{39.47}{0.02} \\
CoTTA    & \res{15.61}{0.00} & \res{35.17}{0.00} & \res{69.44}{0.00} & \res{27.88}{0.00} & \res{37.02}{0.00} \\
LCoTTA   & \res{17.02}{0.94} & \res{34.34}{0.41} & \res{69.36}{0.02} & \res{37.96}{0.96} & \res{39.67}{0.14} \\
AdaDEM   & \ulres{18.59}{1.64} & \res{44.39}{0.23} & \res{68.14}{0.17} & \res{4.96}{0.32} & \res{34.02}{0.37} \\
FOA      & \res{15.49}{0.10} & \res{0.11}{0.00} & \res{69.63}{0.04} & \res{28.88}{0.04} & \res{28.53}{0.03} \\
DeYO     & \res{18.46}{0.80} & \ulres{44.52}{0.13} & \res{69.67}{0.17} & \ulres{40.19}{0.05} & \ulres{43.21}{0.22} \\
\rowcolor{gray!10}
\method{} & \best{21.74}{1.06} & \best{45.25}{0.29} & \best{70.50}{0.16} & \best{41.19}{0.20} & \best{44.67}{0.36} \\
\bottomrule
\end{tabular}
}
\end{minipage}
\hfill
\begin{minipage}[t]{0.39\textwidth}
\centering
\footnotesize
\setlength{\tabcolsep}{4pt}
\captionof{table}{Wild ImageNet-C on ViT-B/16 under label shift, mixed shifts, and BS=1. Unsupported entries are excluded from best-baseline margins.}
% ; collapse entries are shown only as raw official-scope reproductions.}
\label{tab:l4}

\adjustbox{max width=\linewidth}{%
\begin{tabular}{lccc}
\toprule
Method & Label shift & Mixed shifts & BS=1 \\
\midrule
Source & \res{39.75}{0.00} & \res{39.83}{0.00} & \res{39.83}{0.00} \\
Tent & \res{50.33}{0.16} & \res{58.80}{0.16} & \res{5.38}{0.04} \\
EATA & \res{33.89}{0.23} & \res{60.91}{0.23} & \res{2.61}{0.20} \\
SAR & \res{56.52}{0.15} & \res{53.60}{0.15} & \ulres{55.82}{0.47} \\
CoTTA & \res{41.84}{0.01} & \res{42.60}{0.01} & \res{44.35}{0.02} \\
LCoTTA & \res{45.33}{3.34} & \ulres{61.01}{3.34} & \res{32.08}{22.63} \\
AdaDEM & \res{6.01}{0.24} & \res{60.65}{0.24} & \res{0.18}{0.10} \\
% SURGEON & - & - & - \\
FOA & \res{49.24}{0.04} & \res{49.69}{0.04} & \res{0.64}{0.04} \\
DeYO & \ulres{56.66}{0.28} & \res{53.57}{0.28} & \res{5.51}{0.28} \\
% RoTTA & - & - & - \\
\rowcolor{gray!10}
\method{} & \best{57.47}{0.41} & \best{61.67}{0.41} & \best{59.16}{0.41} \\
\bottomrule
\end{tabular}%
}
\end{minipage}
\end{table*}


Natural shifts test transfer to semantic changes rather than synthetic corruptions. Table~\ref{tab:l3} reports ViT-B/16 results on ImageNet-A, ImageNet-R, ImageNet-V2, and ImageNet-Sketch. \method{} reaches $44.67\pm0.36$ average accuracy, $1.46$ points above DeYO. The gains are $3.15$ on ImageNet-A and $0.73$, $0.73$, and $1.00$ on the other three datasets. Because these labels are not induced by one corruption operator, the result supports the portability of \core{} rather than a benchmark-specific bias.

\subsection{Level 3: Wild scenarios and batch-size-one streams}

Wild streams remove common TTA stabilizers and expose failure cases of online self-supervision. Table~\ref{tab:l4} evaluates ImageNet-C wild scenarios under label shift, mixed shifts, and batch-size-one inference on ViT-B/16. \method{} reaches $57.47\pm0.41$, $61.67\pm0.41$, and $59.16\pm0.41$, exceeding the strongest stable compatible baselines by $0.81$, $0.66$, and $3.34$ points.

The batch-size-one column is the clearest stress case because large-batch statistics disappear. There \method{} exceeds SAR, the strongest stable compatible baseline, without relying on unsupported or anomalous entries. This clarifies the role split: \core{} decides whether an entity updates the model, and \care{} bounds the semantic move after acceptance.
Appendix~\ref{app:reproduction_notes} documents anomalous or unsupported baselines, their compatibility scope, default hyperparameters, and failure reasons. Unsupported entries are excluded from best-baseline margins, and collapse entries are reported only as raw official-scope reproductions.

\subsection{Level 4: Dense prediction}
\begin{table*}[t]
\centering
\scriptsize
\setlength{\tabcolsep}{3pt}
\renewcommand{\arraystretch}{1.08}
\caption{Online error (\%) and mIoU for TTA on Cityscapes-to-ACDC with SegFormer-B5. Best in \textbf{bold}; second-best \underline{underlined}.}
\label{tab:timestamp_comparison}
\resizebox{\textwidth}{!}{
\begin{tabular}{c|c|cccc|cccc|cccc|cccc|c|c}
\toprule
\multirow{3}{*}{Architectures}
& \multirow{3}{*}{Methods}
& \multicolumn{16}{l|}{\hspace{1.7em} Time Stamp $\xrightarrow{\hspace{9cm}}$}
& \multirow{3}{*}{Mean (\%) $\downarrow$}
& \multirow{3}{*}{mIoU (\%) $\uparrow$} \\

&
& \multicolumn{4}{c|}{1}
& \multicolumn{4}{c|}{4}
& \multicolumn{4}{c|}{7}
& \multicolumn{4}{c|}{10}
& & \\

&
& Fog & Night & Rain & Snow
& Fog & Night & Rain & Snow
& Fog & Night & Rain & Snow
& Fog & Night & Rain & Snow
& & \\
\midrule

\multirow{5}{*}{SegFormer-B5}
& Source~\cite{xie2021segformer}
& 30.8 & 59.1 & 40.8 & 42.3
& 30.8 & \underline{59.1} & 40.8 & 42.3
& 30.8 & 59.1 & 40.8 & 42.3
& 30.8 & 59.1 & 40.8 & 42.3
& 43.3 & 56.75 \\

& CoTTA~\cite{wang2022cotta}
& \underline{29.2} & \textbf{58.5} & \textbf{37.1} & \underline{40.5}
& \underline{29.2} & 59.2 & \underline{37.4} & \underline{40.5}
& \underline{29.2} & 59.1 & \underline{37.5} & \underline{40.6}
& \underline{29.3} & 58.9 & \underline{37.5} & \underline{40.6}
& \underline{41.5} & \underline{58.24} \\

& MeCTA~\cite{hong2023mecta}
& 30.5 & 60.2 & \underline{40.5} & 42.1
& 33.3 & 62.8 & 42.3 & 45.1
& 35.4 & 65.4 & 46.7 & 47.8
& 37.2 & 68.3 & 47.8 & 50.9
& 47.3 & 53.69 \\

& \textsc{Surgeon~\cite{ma2025surgeon}}
& 31.8 & 59.8 & 41.1 & 43.2
& 31.5 & 59.4 & 39.6 & 42.1
& 30.7 & \underline{58.9} & 38.7 & 41.5
& 30.1 & \underline{58.1} & 38.1 & 41.1
& 42.9 & 57.61 \\

& \textsc{Atlas}
& \textbf{29.0} & \underline{58.8} & \textbf{37.1} & \textbf{40.1}
& \textbf{28.9} & \textbf{58.5} & \textbf{36.9} & \textbf{39.8}
& \textbf{28.4} & \textbf{57.9} & \textbf{36.5} & \textbf{39.5}
& \textbf{28.1} & \textbf{57.5} & \textbf{36.1} & \textbf{39.2}
& \textbf{40.8} & \textbf{59.3} \\

\bottomrule
\end{tabular}
}
\end{table*}

Semantic segmentation directly tests entity-scale normalization because each input contributes a dense pixel lattice. Table~\ref{tab:timestamp_comparison} reports Cityscapes$\rightarrow$ACDC adaptation with SegFormer-B5. \method{} obtains the lowest mean online error, $40.8\%$, and the highest mIoU, $59.30\%$. The closest baseline is CoTTA with $41.5\%$ online error and $58.24\%$ mIoU, so the gain is $0.7$ online-error points and $1.06$ mIoU points.

This is a moderate stability gain rather than the paper's main margin. Because ACDC contains adverse weather and illumination changes where large regions can dominate the loss, \sane{} ties the update scale to selected pixel evidence rather than raw pixel count.

\subsection{Level 5: Vision-language model adaptation}
\begin{table*}[hbtp]
\centering
\caption{CLIP ViT-B/16 prompt-only VLM adaptation on natural-shift test sets. Top-1 accuracy (\%) is reported. OOD Avg. averages ImageNet-A, ImageNet-V2, ImageNet-R, and ImageNet-Sketch.}
\label{tab:l6}
\vspace{2pt}
\resizebox{\textwidth}{!}{%
\begin{tabular}{l|c c c c|c c}
\specialrule{1.2pt}{0pt}{0pt}
Methods & ImageNet-A & ImageNet-V2 & ImageNet-R & ImageNet-S & OOD Avg. & $\Delta$ \\
\specialrule{0.8pt}{0pt}{0pt}
Zero-Shot     & \res{47.4}{0.00} & \res{60.6}{0.00} & \res{74.1}{0.00} & \res{46.0}{0.00} & \res{57.0}{0.00} & +0.0 \\
TPT           & \res{54.8}{0.08} & \res{62.5}{0.02} & \res{75.5}{0.08} & \res{47.3}{0.20} & \res{60.0}{0.10} & \gain{+3.0} \\
TDA           & \res{49.1}{0.00} & \best{64.6}{0.00} & \res{74.0}{0.00} & \res{47.7}{0.00} & \res{58.9}{0.00} & \gain{+1.9} \\
AdaDEM        & \ulres{55.0}{0.12} & \res{62.0}{0.04} & \ulres{77.9}{0.74} & \ulres{48.0}{0.04} & \ulres{60.7}{0.24} & \gain{\underline{+3.7}} \\
\rowcolor{gray!15}
\method{}     & \best{56.8}{0.73} & \ulres{63.9}{0.08} & \best{78.5}{0.06} & \best{49.1}{0.09} & \best{62.1}{0.24} & \gain{\textbf{+5.1}} \\
\specialrule{0.8pt}{0pt}{0pt}
CoOp          & \res{49.8}{0.00} & \res{64.6}{0.00} & \res{74.8}{0.00} & \res{47.2}{0.00} & \res{59.1}{0.00} & \gain{+2.1} \\
TPT           & \res{59.1}{0.74} & \res{64.5}{0.05} & \ulres{78.6}{0.01} & \ulres{49.6}{0.09} & \res{63.0}{0.22} & \gain{+6.0} \\
TDA           & \res{56.0}{0.39} & \res{63.5}{0.00} & \res{76.3}{0.06} & \res{46.4}{0.07} & \res{60.6}{0.13} & \gain{+3.6} \\
AdaDEM        & \ulres{61.8}{0.72} & \ulres{64.8}{0.76} & \res{77.0}{0.02} & \res{49.3}{0.00} & \ulres{63.2}{0.38} & \gain{\underline{+6.2}} \\
\rowcolor{gray!15}
\method{}     & \best{62.5}{0.46} & \best{65.6}{0.01} & \best{78.8}{0.85} & \best{49.9}{0.06} & \best{64.2}{0.35} & \gain{\textbf{+7.2}} \\
\specialrule{1.2pt}{0pt}{0pt}
\end{tabular}%
}
\end{table*}
\begin{table}[t]
\centering
\caption{Protocol metadata for the VLM comparison. ``Reset'' means prompt and optimizer state are restored for each test instance. Source ensemble is used only where the method explicitly requests it.}
\label{tab:vlm_protocol}
\small
\resizebox{\columnwidth}{!}{%
\begin{tabular}{lccccc}
\toprule
Method & Runtime & Reset & Trainable parameters & Views & Source ensemble \\
\midrule
Source & episodic eval. & -- & none & original & optional prompt templates \\
TPT & episodic & yes & prompt context & augmented & no \\
AdaDEM & episodic & yes & prompt context & augmented & no \\
ATLAS-Core & episodic & yes & prompt context & selected augmented & yes \\
ATLAS-Full & episodic & yes & prompt context & original + augmented & yes \\
TDA & online & no & cache state & original & no \\
\bottomrule
\end{tabular}}
\end{table}


Prompt-only VLM adaptation asks whether evidence selection, anchoring, and selected-entity normalization remain useful when the adapted object is a prompt-view prediction rather than a normalization layer. VLM experiments use episodic instance-wise prompt adaptation: the prompt context and optimizer state are reset to their initialization for every test instance. This is not the continual stream protocol used for classification. Table~\ref{tab:l6} reports CLIP ViT-B/16 results under zero-shot and CoOp prompt initializations, while Table~\ref{tab:vlm_protocol} makes the protocol differences explicit. The historical Full results are retained pending the registered Core/Full rerun; until that gate completes, VLM is treated as an extension rather than evidence for an identical cross-task loss.

\subsection{Ablation and same-role replacements}
\begin{table*}[t]
\centering
\small
\caption{Orthogonal decomposition and same-role replacements on ImageNet-C with ViT-B/16. A0 keeps the shared online optimizer and parameter-update scope but replaces \core{} and \care{} with raw entropy minimization over all samples; A1 applies the same entropy loss only to \core{}-selected entities. Replacement rows compare \core{} and \care{} with published alternatives of similar role.}
\label{tab:l7}
\resizebox{\textwidth}{!}{%
\begin{tabular}{lcccl}
\toprule
Config & \core{} & \care{} & Continual mean acc. & Comment \\
\midrule
A0 & None & None & \res{15.58}{6.73} & raw entropy on all samples \\
A1 & \core{} & None & \res{34.86}{11.67} & raw entropy on selected entities \\
A2 & None & \care{} & \res{4.90}{0.20} & conservative anchoring only \\
A3 & \core{} & \care{} & \res{65.14}{0.05} & full \method{} \\
B1 & DeYO-style selection~\cite{lee2024deyo} & \care{} & \res{61.4}{0.13} & replace \core{} with a published selector \\
B2 & \core{} & AdaDEM-style decoupled entropy~\cite{ma2025adadem} & \res{62.8}{0.21} & replace \care{} with a published decoupled objective \\
\bottomrule
\end{tabular}}
\end{table*}
\begin{wraptable}{r}{0.50\textwidth}
\vspace{-10pt}
\centering
\scriptsize
\setlength{\tabcolsep}{4pt}
\renewcommand{\arraystretch}{1.06}
\caption{SegFormer-B5 SANE ablation under the same Cityscapes$\rightarrow$ACDC continual protocol as Table~\ref{tab:timestamp_comparison}. \textsc{Atlas w/o SANE} replaces selected-pixel normalization with total-pixel normalization.}
\label{tab:sane_seg_placeholder}
\begin{tabular}{lcc}
\toprule
Variant & Mean online error (\%) $\downarrow$ & mIoU (\%) $\uparrow$ \\
\midrule
\textsc{Atlas} & \textbf{40.80} & \textbf{59.30} \\
\textsc{Atlas w/o SANE} & 98.92 & 1.06 \\
\bottomrule
\end{tabular}
\vspace{-8pt}
\end{wraptable}


The ablation in Table~\ref{tab:l7} verifies that \core{} and \care{} are complementary on ViT-B/16 continual ImageNet-C, and Appendix~\ref{app:ablation_details} reports the corruption-wise expansion. The base online adapter A0 keeps the shared optimizer and update scope but replaces both \core{} and \care{} with raw entropy minimization over all samples, yielding $15.58\pm6.73$. Adding \core{} alone raises accuracy to $34.86\pm11.67$, while keeping \care{} without \core{} reaches only $4.90\pm0.20$. The full configuration reaches $65.14\pm0.05$. Replacing \core{} with DeYO-style selection~\cite{lee2024deyo} or \care{} with AdaDEM-style decoupled entropy minimization~\cite{ma2025adadem} lowers accuracy to $61.4\pm0.13$ and $62.8\pm0.21$.

Table~\ref{tab:sane_seg_placeholder} shows the SegFormer-B5 SANE ablation under the same continual protocol: replacing selected-pixel normalization with total-pixel normalization raises mean online error from $40.80\%$ to $98.92\%$ and reduces mIoU from $59.30\%$ to $1.06\%$.

These comparisons show that ATLAS depends on corroborative selection, conservative anchoring, and selected-entity normalization together. In particular, \care{} is effective only when paired with \core{}, and dense prediction is highly sensitive to normalization by selected rather than total entities.
